% =====================================================================
%  3. Problema de investigación
% =====================================================================
\section{Problema de investigación}
\label{sec:problema}

\subsection{Delimitación del objeto}
\label{sec:problema-objeto}

El objeto de esta investigación no es el repositorio institucional ni el recurso
educativo abierto, sino el \textbf{corpus documental de un curso}: el conjunto de
guías, apuntes, presentaciones, lecturas, listados de preguntas frecuentes y
evaluaciones que un docente produce, edita y reutiliza a lo largo de varios
semestres. Esta delimitación importa porque desplaza el problema respecto de la
literatura de calidad de repositorios revisada en la
sección~\ref{sec:contexto-calidad}:

\begin{itemize}
  \item La unidad de análisis no es el recurso aislado, sino la
        \textbf{relación entre documentos}. Un documento puede ser individualmente
        correcto y, aun así, contradecir a otro del mismo corpus.
  \item No existe una comunidad de usuarios que ejerza control colaborativo de
        calidad. El aseguramiento recae sobre una persona.
  \item El corpus es \textbf{cerrado y conocido}, lo que hace posible verificar cada
        afirmación contra documentos concretos en lugar de contra el conocimiento
        general del mundo.
\end{itemize}

\subsection{Manifestaciones del problema}
\label{sec:problema-manifestaciones}

Sobre ese corpus se observan cuatro manifestaciones, que esta investigación trata
como clases distintas y no como un único fenómeno de \enquote{desorden}:

\begin{enumerate}
  \item \textbf{Redundancia.} El mismo contenido aparece en varios documentos con
        redacciones distintas. Dispersa la referencia autorizada y multiplica los
        puntos que hay que actualizar cuando algo cambia.
  \item \textbf{Inconsistencia o contradicción.} Dos documentos del corpus afirman
        cosas incompatibles: un valor numérico, una definición, un procedimiento, un
        plazo, un criterio de evaluación.
  \item \textbf{Desactualización.} Un contenido que fue correcto deja de serlo
        porque cambió el estándar, la versión de una tecnología o la normativa, y
        permanece en circulación. Incluye el caso que
        \textcite{leacock2007framework} señalan como degradación de la reusabilidad:
        información específica de una instancia pasada (fechas, nombres, cupos).
  \item \textbf{Ausencia de evidencia localizable.} No hay un punto desde el cual
        consultar el corpus y establecer de qué documento y de qué fragmento
        proviene una afirmación.
\end{enumerate}

La distinción entre las clases 2 y 4 no es retórica. La literatura de inferencia en
lenguaje natural la trata como una distinción de clase: \textcite{koreeda2021contractnli}
en inferencia a nivel documental sobre contratos, \textcite{thorne2018fever} en
extracción y verificación de hechos, y \textcite{gubelmann2024capturing} en su
revisión de conjuntos de datos de NLI separan \emph{contradice} de
\emph{no hay información suficiente}. En estas tareas la contradicción es la clase
más difícil y la clase neutral la peor reconocida por los modelos. Un sistema que
colapse ambas categorías en \enquote{alerta} produce un tipo de error que el docente
no puede interpretar: no distingue \enquote{esto se contradice con aquello} de
\enquote{no encontré con qué comparar esto}.

\subsection{La tensión que constituye el problema de investigación}
\label{sec:problema-tension}

Asistir la curación con modelos de lenguaje introduce el riesgo que la propia tarea
pretende reducir. \textcite{huang2025survey} distinguen la alucinación de
\emph{factualidad} --- contradecir o inventar hechos del mundo --- de la de
\emph{fidelidad} --- apartarse del contexto proporcionado o contradecirlo --- y
documentan que la recuperación aumentada no resuelve el problema:
\textcquote[p.~26]{huang2025survey}{Despite being designed to mitigate LLM
hallucinations, retrieval-augmented LLMs can still produce hallucinations}.

El punto crítico para este proyecto está en las causas que identifican. Entre las
fuentes de alucinación en sistemas RAG figuran las \textbf{fuentes con información
incorrecta o desactualizada}, la fragmentación inadecuada del texto, el contexto
ruidoso y los conflictos de conocimiento. En un corpus de curso con contradicciones,
\emph{los propios documentos son esa fuente de riesgo}: el sistema recupera evidencia
inconsistente y puede fundamentar sobre ella una sugerencia incorrecta. A ello se
suma que los modelos no reconocen bien sus límites de conocimiento y producen
falsedades con confianza, lo que vuelve insuficiente cualquier señal de confianza que
no haya sido contrastada con la tasa real de acierto.

De esta tensión surge el problema:

\begin{estado}[colframe=ecAzul]
\small
\textbf{Problema de investigación.} Un sistema que asista la curación de un corpus
de curso debe operar sobre un material que puede ser internamente inconsistente,
usando un componente generativo que puede producir afirmaciones no fundamentadas y
señales de confianza que pueden no corresponder a su tasa de acierto. Se desconoce
\emph{en qué medida} una arquitectura que combina recuperación, detectores,
razonamiento con herramientas, evidencia estructurada y revisión humana detecta las
inconsistencias reales de ese corpus, \emph{qué componentes} de esa combinación
explican su comportamiento, y \emph{con qué efecto} sobre el trabajo del docente.
\end{estado}

\subsection{Lo que el problema no incluye}
\label{sec:problema-exclusiones}

Para evitar que el objeto se desdibuje, se excluyen explícitamente:

\begin{itemize}
  \item La \textbf{verificación factual contra el mundo}. El criterio de corrección
        es la consistencia con el corpus del curso y, cuando aplique, con las fuentes
        externas que el docente autorice, no la verdad general de una afirmación.
  \item La \textbf{evaluación del aprendizaje}. No se mide si un corpus más
        coherente mejora resultados de aprendizaje; esa relación excede el alcance y
        requiere un diseño distinto.
  \item La \textbf{calidad pedagógica} en las dimensiones de LORI ajenas al
        contenido (motivación, diseño de la presentación, accesibilidad). El sistema
        no juzga si un material es didácticamente adecuado.
  \item La \textbf{sustitución del juicio docente}. La decisión sobre qué se
        conserva, qué se corrige y qué se descarta permanece fuera del sistema por
        diseño, no por limitación técnica.
\end{itemize}
